% ==================================================================
%  徐森林、薛春华《数学分析》（第一册）
%  第 4 章复习题
% ==================================================================
%  源：原书印刷页 279-281（PDF 物理页 285-287）
%  底本：PaddleOCR-VL-1.6 + 高清影印逐题核验校订
%  规范：exblock 容器; 全西文半角标点; 单一落盘出口 _texout.py
% ==================================================================
\begin{exblock}{复习题 4}

\item $\mathbb{R}$ 上的 $2$ 阶可导函数 $f(x)$ 满足 $f(0)=f'(0)=0$, 且 $|f''(x)|\leqslant C|f(x)f'(x)|$, $\forall x\in\mathbb{R}$, 其中 $C$ 为正的常数. 证明: $f(x)\equiv 0$.

\item 设 $\mathbb{R}$ 上 $n$ 阶可导函数 $f(x)$ 满足 $f(0)=f'(0)=\cdots=f^{(n-1)}(0)=0$, 且存在正常数 $C$ 与固定的 $j\in\{0,1,\cdots,n-1\}$, 使 $|f^{(n)}(x)|\leqslant C|f^{(j)}(x)|$, $\forall x\in\mathbb{R}$ ($f^{(0)}(x)=f(x)$). 证明: $f(x)\equiv 0$, $\forall x\in\mathbb{R}$.

\item 设 $f$ 在 $(a,b)$ 内无穷阶可导, 且各阶导数均只取正值. 证明: 对 $\forall x_0 \in (a,b), \exists r>0$, s.t. 当 $x\in[x_0-r, x_0+r]\subset(a,b)$ 时, 有
\[
f(x)=\sum_{k=0}^{\infty}\frac{f^{(k)}(x_0)}{k!}(x-x_0)^k.
\]

\item 设函数 $f$ 在 $[0,2]$ 上满足 $|f(x)|\leqslant 1$ 及 $|f''(x)|\leqslant 1$. 证明: 在区间 $[0,2]$ 上 $f'(x)$ 有界, 且 $2$ 是最小的界.

\item 设 $P_n(x)=1+\frac{x}{1!}+\frac{x^2}{2!}+\cdots+\frac{x^n}{n!}, n\in\mathbb{N}$.
\begin{enumerate}[label=(\arabic*)]
  \item 当 $n$ 为偶数时, $P_n>0$;
  \item 当 $n$ 为奇数时, $P_n$ 有惟一的实零点;
  \item $P_{2n+1}$ 的实零点记为 $x_n$ ($n=0,1,2,\cdots$).
\end{enumerate}
证明: 数列 $x_n$ 严格单调减且趋于 $-\infty$.

\item 证明: 多项式
\[
\sum_{k=1}^{n}\frac{(2x-x^2)^k-2x^k}{k}
\]
能被 $x^{n+1}$ 整除. (提示: 对 $-\ln(1-x)^2=-2\ln(1-x)$ 与 $\ln(1-x)^2=\ln(1-2x+x^2)$ 作带 Peano 型余项的 Taylor 展开.)

\item 设 $f(x)=\lim_{n\to+\infty}n^x\left[\left(1+\frac{1}{n+1}\right)^{n+1}-\left(1+\frac{1}{n}\right)^n\right]$, 求 $f(x)$ 的定义域与值域.

\item 设 $f$ 为 $[a,b]$ 上的 $2$ 阶可导函数, 且满足
\[
[f(x)]^2+[f''(x)]^2=r^2\quad (r>0).
\]
证明: $|f'(x)|\leqslant\left(\frac{2}{b-a}+\frac{b-a}{2}\right)r, \forall x\in[a,b]$.

\item 设 $f$ 在 $[0,+\infty)$ 上 $2$ 阶可导, $f''(x)$ 有界, $\lim_{x\to+\infty}f(x)=0$. 证明: $\lim_{x\to+\infty}f'(x)=0$.

\item 设 $f$ 在 $\mathbb{R}$ 上有各阶导数, 且 $\exists M>0$, s.t. 对 $\forall n\in\{0,1,2,\cdots\}$, 有 $|f^{(n)}(x)|\leqslant M, \forall x\in\mathbb{R}$. 如果在一个无限有界集 $E$ 上 $f=0$, 证明: 在 $\mathbb{R}$ 上 $f=0$.

\item 设 $f(x)$ 在闭区间 $[a,b]$ 上 $2$ 阶可导, $f'\left(\frac{a+b}{2}\right)=0$.
\begin{enumerate}[label=(\arabic*)]
  \item 证明: $\exists\xi\in(a,b)$, s.t.
  \[
  |f''(\xi)|\geqslant\frac{4}{(b-a)^2}|f(b)-f(a)|;
\]
  \item 说明常数 $4$ 是最好的. 即对 $\forall M>4$, 总可找到一个具体的区间 $[a,b]$ 及满足条件的 $f(x)$, 使对 $\forall\xi\in(a,b)$, 都有
  \[
  |f''(\xi)|<\frac{4}{(b-a)^2}|f(b)-f(a)|;
\]
  \item 如果再设 $f(x)\not\equiv \text{常数}$, 则 $\exists\xi\in(a,b)$, s.t.
  \[
  |f''(\xi)|>\frac{4}{(b-a)^2}|f(b)-f(a)|.
\]
\end{enumerate}

\end{exblock}
